\documentclass[letterpaper]{article}
\usepackage[draft]{aaai2026}
\usepackage{times}
\usepackage{helvet}
\usepackage{courier}
\usepackage[hyphens]{url}
\usepackage{graphicx}
\urlstyle{rm}
\def\UrlFont{\rm}
\usepackage{natbib}
\usepackage{caption}
\usepackage{amsmath}
\frenchspacing
\setlength{\pdfpagewidth}{8.5in}
\setlength{\pdfpageheight}{11in}
\ifdefined\pdfinfo
\pdfinfo{/TemplateVersion (2026.1)}
\fi
\setcounter{secnumdepth}{0}
\title{Ranking-Aware Concept-Level Sparse Retrieval\\with Sparse Autoencoders}
\author{Team Author List Pending}
\affiliations{Nanyang Technological University\\AI6103 Course Project}
\newcommand{\pending}[1]{\par\noindent\textbf{Draft---#1}\par}
\begin{document}
\maketitle
\begin{abstract}
Sparse autoencoders can expose interpretable latent concepts in dense retrieval
embeddings and support retrieval through an inverted index. However, embedding
reconstruction does not directly optimize passage ranking, and the released
concept-level sparse retrieval pipeline uses a fixed BM25-style scoring rule.
This project investigates modifications to scoring, SAE training, passage
sparsity, concept pruning, and query expansion. We report a completed scoring
experiment that freezes the embeddings, SAE, and passage postings and learns
additive concept impacts by distillation from SimLM. Over 8.84 million passages,
three seeds achieve MS MARCO dev MRR@10 of $0.3600\pm0.0015$, compared with
$0.3239$ for validation-tuned concept BM25. Retrieval is slower and TREC
Recall@1000 decreases, despite improved top-10 ranking.
\textbf{Draft: team results and combined conclusions remain to be integrated.}
\end{abstract}
% Owner: P1. Keep the final introduction around 0.6 pages.
\section{Introduction}
Passage retrieval selects evidence for question answering from a large corpus.
Dense retrievers encode questions and passages into continuous vectors, enabling
semantic matching beyond exact word overlap. These representations can retrieve
effectively, but their individual dimensions are difficult to interpret, and
storing a dense vector for every passage is expensive.

\citet{park-etal-2025-decoding} introduce sparse autoencoders (SAEs) to represent
dense retrieval embeddings through a small set of active latent concepts.
Concept-Level Sparse Retrieval (CL-SR) searches these representations with an
inverted index. This exposes individual matching concepts and offers a route
to reducing storage. Two objectives remain distinct: reconstructing an embedding
and ranking relevant passages. A low reconstruction error alone does not ensure
that the final concept scores preserve the teacher's retrieval preferences.

Our project asks whether explicitly optimizing ranking, together with controlling
concept usage, improves this quality--efficiency trade-off. We investigate five
stages independently: learned impact scoring (P1), ranking-aware SAE training
(P2), adaptive passage sparsity (P3), concept pruning (P4), and query concept
expansion (P5). Each experiment must use a measured baseline under its own
controlled protocol. Combining results across stages additionally requires
compatible checkpoints, corpus statistics, and evaluation settings.

P1 isolates the scoring rule. A small neural network learns each shared concept's
contribution from dense SimLM scores, while keeping the indexed representations
fixed. Its additive form permits full-corpus concept retrieval. We evaluate
against released and validation-tuned BM25, with held-out queries excluded from
selection. The completed experiment improves top-10 quality, but also exposes
costs in latency and deeper TREC recall. Reporting these costs is necessary to
assess whether a ranking gain is practically useful.

% Owner: P2. Target about 0.5 pages; add verified BibTeX entries.
\section{Related Work}
The starting point is the SAE-based interpretation and discretization of dense
retrieval embeddings by \citet{park-etal-2025-decoding}. Our scoring contribution
keeps that representation fixed and learns how shared activations contribute
to ranking.
\pending{P2: synthesize dense and learned sparse retrieval, SAE retrieval,
SPLARE, and the NAACL 2025 baseline. Verify the exact prior objective and
bibliographic metadata before contrasting it with P2.}

% Owner: P3. Target about 0.5 pages; verify against released code.
\section{Base Retrieval Pipeline}
SimLM maps a query or passage to a 768-dimensional embedding. A frozen SAE
encodes each embedding into sparse concept activations. Passage activations
are stored in posting lists indexed by concept; a query visits lists for its
active concepts. The final score sums contributions from concepts shared by
the query and passage. This system retrieves passages; answer generation is
outside the evaluated pipeline.

For P1, the released BM25-style rule is reproduced with signed concept IDF,
document length equal to retained concept count, and constants
$k_1=0.6$, $k_2=1.5$, and $b=2.0$. The tuned baseline changes $k_1$ and $b$
using validation only. Every P1 method uses the same passage postings and
corpus statistics.
\pending{P3: complete the SAE training objective, BatchTopK and inference
threshold explanation. Distinguish training k=32 from the P1 passage cap of 24.}

\section{Methods}
% Owner: P1. Verified against p1/model.py and the frozen run protocol.
\subsection{P1: Learned Impact Scoring}
For query $q$, passage $d$, and a shared concept $c$, let $q_c$ and $d_c$ denote
the activation magnitudes. Define the feature vector
\begin{equation}
\begin{split}
\phi_c(q,d)=\big[&\log(1+q_c),\log(1+d_c),\\
&\mathrm{IDF}_c,\log(1+\ell_d/\bar\ell)\big].
\end{split}
\end{equation}
where $\ell_d$ is the retained passage concept count, $\bar\ell$ its corpus
mean, and $\mathrm{IDF}_c=\log[(N-\mathrm{df}_c+0.5)/(\mathrm{df}_c+0.5)]$.
The learned score is
\begin{equation}
s_\theta(q,d)=\sum_{c\in\mathrm{supp}(q)\cap\mathrm{supp}(d)}
f_\theta(\phi_c(q,d)).
\end{equation}
The MLP $f_\theta$ has four inputs, 32 ReLU hidden units, and one softplus
output, totaling 193 trainable parameters. The nonnegative output defines
each shared concept's contribution. The dense teacher is used for training;
inference scores postings over the full corpus without a dense candidate filter.

For each training query, SimLM provides scores $t_i$ for 100 dense candidates.
We minimize $\mathrm{KL}(p_T\|p_\theta)$, where
$p_T(i)=\mathrm{softmax}(t)_i$ and
$p_\theta(i)=\mathrm{softmax}(s_\theta)_i$, at temperature one.
Embeddings, SAE parameters, and passage postings remain fixed. Adam uses
learning rate $10^{-3}$ for ten epochs. Seeds 42, 43, and 44 use the same
5,000 training queries; each checkpoint is selected by validation MRR@10.

% Contributors edit their own method subsection here or add a separate input.
\subsection{P2: Ranking-Aware SAE Training}
\pending{P2: final loss, controls, validation-selected settings, NMSE constraint,
and any deviations from the original plan.}
\subsection{P3: Adaptive Passage Sparsity}
\pending{P3: threshold rule with minimum and maximum concept counts, fixed-k
baselines, and selection protocol.}
\subsection{P4: Concept Pruning}
\pending{P4: learned gate, regularization, DF baseline, and impact-index
quantization. State how negative IDF contributions are handled.}
\subsection{P5: Query Concept Expansion}
\pending{P5: expansion network, teacher objective, decoder-similarity baseline,
and validation-only selection.}

% Owner: P4. The paragraphs below describe P1's verified setup only.
\section{Experimental Setup}
\paragraph{P1 data and selection.}
We index all 8,841,823 MS MARCO passages using the shared
\texttt{B0\_final\_atom} SAE and \texttt{data\_titled} artifacts. Its stored
inference threshold is 0.49723; passages retain at most 24 concepts.
Training uses 5,000 sampled queries and a separate shared 5,000-query validation
split. Released and learned scoring share the identical index.

The BM25 grid has 28 settings: $k_1\in\{0.3,0.6,1.2,2.4\}$ and
$b\in\{0,0.25,0.5,0.75,1,1.5,2\}$. Screening uses 500 validation queries;
the best three settings, the released setting, and the previous candidate-tuned
setting are confirmed on all 5,000. The selected setting is
$k_1=1.2$, $b=1.0$, $k_2=1.5$. All settings and checkpoint hashes were frozen
before dev or TREC evaluation.

\paragraph{Metrics and inference.}
MS MARCO dev has 6,980 judged queries; MRR@10 is primary. TREC DL 2019 and
2020 contribute 43 and 54 judged queries, evaluated with linear-gain NDCG@10.
Recall@1000, serialized index bytes, and query latency measure complementary
quality and efficiency. Neural results report three-seed means and sample
standard deviations. The predeclared primary seed 42 is compared with both
baselines using 10,000 paired randomizations, with Holm correction across
the two comparisons within each split. TREC tests are secondary.

\paragraph{Audit and timing.}
Training and validation query IDs are disjoint. A text audit found two
case/punctuation-normalized train--validation matches; excluding them did not
explain the validation gain. No corresponding train-text matches were found
among judged held-out queries. These checks do not prove that overfitting is
impossible. Warm retrieval latency is measured on an NVIDIA A40, including
posting transfers, scoring, and ranking, but excluding text/SAE encoding and
index loading. Float64 accumulation and numeric passage-ID tie breaking are
shared across P1 scoring methods.
\pending{P4: integrate the shared setup and identify method-specific
checkpoint, index-format, and hardware differences before comparing rows.}

\section{Results and Analysis}
\subsection{P1: Held-Out Full-Corpus Retrieval}
Table~\ref{tab:p1-quality} compares scoring methods with identical postings.
P1 improves dev MRR@10 by 0.0361 (11.1\% relative) over tuned BM25.
For the predeclared seed 42, improvements over both baselines are significant
on dev (Holm-adjusted $p=0.0002$). DL20 improves significantly ($p=0.0034$
versus tuned BM25); DL19 does not ($p=0.1262$). We do not select a seed on test results.
\begin{table}[t]
\centering
\small
\begin{tabular}{lrrr}
\hline
Metric & Released & Tuned & P1 \\
\hline
Dev MRR@10 & 0.3158 & 0.3239 & $0.3600 \pm 0.0015$ \\
DL19 NDCG & 0.6111 & 0.6166 & $0.6502 \pm 0.0047$ \\
DL20 NDCG & 0.5685 & 0.5727 & $0.6195 \pm 0.0056$ \\
\hline
\end{tabular}
\caption{Full-corpus held-out quality (NDCG@10 on TREC). P1 is mean $\pm$ sample standard deviation over three seeds.}
\label{tab:p1-quality}
\end{table}
\paragraph{Coverage and efficiency.}
Dev Recall@1000 increases from 0.9463 for tuned BM25 to 0.9520 for P1.
On DL19 and DL20 it falls from 0.7086/0.7023 to 0.6677/0.6783.
The shared index occupies 1,688,270,708 bytes (1.572 GiB), with 23.61 concepts
per passage on average. P1 adds a 152,329-byte primary-seed checkpoint.
Mean dev warm retrieval latency is 191.81 ms across P1 seeds versus 122.09 ms
for tuned BM25, approximately 57\% higher. Scoring improves top-10 quality
while preserving postings; it does not improve speed.
\paragraph{Vocabulary mismatch.}
P5 supplied dev subsets derived from lexical BM25 at $k_1=0.9$, $b=0.4$.
Table~\ref{tab:p1-mismatch} evaluates frozen full-corpus rankings on those IDs.
MRR@10 and Recall@1000 improve on all subsets. We verified counts, nesting,
split membership, and ranking checksums. The source lexical run and explicit
generation rule were not supplied, so construction was not independently
reproduced. These are secondary descriptive results without further tuning.
\begin{table}[t]
\centering
\small
\begin{tabular}{lrrr}
\hline
Subset & Queries & Tuned MRR & P1 MRR \\
\hline
Top-10 & 4289 & 0.2018 & 0.2408 \\
Top-100 & 2327 & 0.1385 & 0.1746 \\
Top-1000 & 988 & 0.1087 & 0.1340 \\
\hline
\end{tabular}
\caption{Supplied dev mismatch subsets. P1 values are three-seed means. Recall@1000 improves from 0.9188/0.8691/0.7677 to 0.9275/0.8847/0.7985.}
\label{tab:p1-mismatch}
\end{table}

\subsection{Other Contributions and Combined Ablation}
\pending{Each owner: insert verified results, including negative findings.
P5: assemble the combined ablation only after confirming comparable artifacts
and settings. P3 currently reports full-corpus validation results, not held-out
dev/TREC results. P4 uses a different impact-index backend, so its latency and
size cannot be interpreted as scoring-only effects relative to P1.}

% Owner: P5. Replace the provisional conclusion after team integration.
\section{Conclusion}
The completed P1 experiment shows that a small learned additive scorer can
improve full-corpus top-10 retrieval over released and tuned concept BM25 while
keeping passage postings fixed. The gains extend to supplied vocabulary-mismatch
subsets, but incur higher latency and lower TREC Recall@1000. Thus, better
top-10 ranking does not imply uniformly better retrieval efficiency or coverage.
\pending{P5: conclude across all five measured contributions and combinations;
do not claim unmeasured improvements.}

\bibliography{references}
\end{document}
